\phantomsection
\section*{РЕФЕРАТ}

% TODO: уточнить число страниц/рисунков/таблиц после готовности введения и главы 1.
Выпускная квалификационная работа по теме <<\varTheme>> содержит \_\_\_ страниц печатного текста, \_\_\_ рисунков, \_\_\_ таблиц, 19 использованных источников.

РОТОР, МНОГОСТУПЕНЧАТЫЙ ЦЕНТРОБЕЖНЫЙ НАСОС, ПОДШИПНИК СКОЛЬЖЕНИЯ, ЩЕЛЕВОЕ УПЛОТНЕНИЕ, КРИТИЧЕСКАЯ СКОРОСТЬ, ДИНАМИЧЕСКИЕ КОЭФФИЦИЕНТЫ, УСТОЙЧИВОСТЬ РОТОРА, ЛОГАРИФМИЧЕСКИЙ ДЕКРЕМЕНТ, ВИБРАЦИЯ.

Актуальность работы обусловлена широким применением многоступенчатых секционных центробежных насосов типа ЦНС в нефтегазовой отрасли и высокими требованиями к их вибрационной надёжности. Динамическое состояние ротора определяется совместным влиянием гидродинамических подшипников скольжения и щелевых уплотнений, поэтому корректный учёт этих факторов при расчёте критических скоростей, орбит и запаса устойчивости является необходимым условием безотказной эксплуатации.

Объект исследования -- ротор многоступенчатого центробежного насоса ЦНС 105-392 (5МС-10) с гидродинамическими подшипниками скольжения и щелевыми уплотнениями.

Цель работы -- определение и анализ динамических характеристик ротора насоса ЦНС 105-392 с учётом совместного влияния подшипников скольжения и щелевых уплотнений на основе гидродинамического анализа.

Для достижения цели решены следующие задачи:
\begin{itemize}
    \item построена расчётная схема ротора с сосредоточенными массами рабочих колёс, радиальными опорами и щелевыми уплотнениями;
    \item выполнен расчёт поля давления в подшипнике скольжения по уравнению Рейнольдса и определены динамические коэффициенты жёсткости и демпфирования;
    \item собрана двумерная балочная модель ротора, определены собственные частоты, формы изгиба и критические скорости;
    \item рассчитаны вынужденные орбиты, перекос шейки и выполнен анализ устойчивости по логарифмическому декременту;
    \item проведена нормативная оценка вибрационного состояния и параметрические исследования чувствительности результатов.
\end{itemize}

В результате работы определены критические скорости ротора, оценён запас устойчивости и вибрационное состояние для трёх эксплуатационных сценариев, а также выявлены параметры, наиболее существенно влияющие на динамические характеристики ротора.
